\documentclass{article}

\usepackage[T1]{fontenc}
\usepackage{microtype}
\usepackage{graphicx}
\usepackage{booktabs}
\usepackage{multirow}
\usepackage{amsmath}
\usepackage{amssymb}
\usepackage{mathtools}
\usepackage{xcolor}
\usepackage{hyperref}
\usepackage{url}
\usepackage{enumitem}
\usepackage{algorithm}
\usepackage{algorithmic}
\usepackage{cleveref}
\usepackage{balance}

% Use the supplied ICML 2026 two-column style in anonymous review mode.
\usepackage{icml2026}

% Identify the review venue without changing the supplied typesetting contract.
\makeatletter
\renewcommand{\Notice@String}{Preliminary work. Under review for a NeurIPS workshop. Do not distribute.}
\makeatother
\hypersetup{pdfsubject={Anonymous NeurIPS workshop submission}}

\definecolor{resultred}{HTML}{B93835}
\definecolor{resultblue}{HTML}{2468A2}
\newcommand{\trtwo}{\textsc{TR2-Color}}
\newcommand{\trthree}{\textsc{TR3-Color}}
\newcommand{\trfour}{\textsc{TR4-Color}}
\newcommand{\lewm}{\textsc{LeWM}}
\newcommand{\pldm}{\textsc{PLDM}}
\newcommand{\lpwm}{\textsc{LpWM}}
\newcommand{\dinowm}{\textsc{DINO-WM}}
\newcommand{\experimentneeded}[1]{\textcolor{resultred}{\textbf{[Experiment needed:} #1\textbf{]}}}
\newcommand{\artifactneeded}[1]{\textcolor{resultred}{\textbf{[Artifact needed:} #1\textbf{]}}}

\icmltitlerunning{When LeWM Is Color Blind}

\begin{document}

\twocolumn[
\icmltitle{When LeWM Is Color Blind:\\A Two-Room Color-Identity Failure}

  \begin{icmlauthorlist}
    \icmlauthor{Anonymous Author(s)}{anon}
  \end{icmlauthorlist}
  \icmlaffiliation{anon}{Anonymous Institution}
  \icmlcorrespondingauthor{Anonymous Author}{anonymous@domain.com}
  \icmlkeywords{world models, joint embeddings, representation learning, color}
  \vskip 0.25in
]

\printAffiliationsAndNotice{}

\begin{abstract}
We introduce \emph{Tunnel Rooms Color}, a diagnostic environment in which
geometrically identical portals appear at randomized locations and color
determines a delayed teleport destination.  The task separates access to
portal geometry from color--location binding.  Across three two-room training
seeds, LeWorldModel (LeWM) spatial states localize position-defined hubs to
$6.41\pm0.34$ pixels, yet color-conditioned error remains $25.97\pm0.20$
pixels against a 30.02-pixel bias-only reference; CLS and MLP states are weaker.
No tested LeWM global readout is strong, whereas PLDM spatial states reach
$18.05\pm0.41$ pixels and LpWM CLS states $15.65\pm1.42$.  The failure is
stage-specific: native LeWM patches retain fixed-slot hue at 96--100\%
balanced accuracy, but all 12 final projectors in our inverse-dynamics
factorial remain near chance (45.0--56.3\%).  At selected spatial checkpoints,
PLDM and LpWM projectors retain visible hue at 100\%; only LpWM also retains
destination binding.  Detached MSE decoders retain 4.85\% of original
green--orange separation for LeWM, versus 49.5\% for PLDM and 37.8\% for LpWM;
LeWM instead reconstructs their olive average and repeats the behavior for
green/blue.  Inverse-dynamics and temporal-alignment sweeps do not rescue the
two-room failure, while spatial three-room LeWM is a counterexample.  Thus the
evidence identifies a consistent loss of accessible TR2 color along LeWM's
learned global path, not an absence of color in its backbone patches or
universal color sensitivity in every comparator state.
\end{abstract}

\begin{figure*}[t]
    \centering
    \includegraphics[width=\textwidth]{figures/colorblind_teaser.pdf}
    \caption{\textbf{LeWM averages two predictive colors, although selected
    comparators and three-room LeWM can preserve separation.}  Top: on one
    matched final-epoch TR2 frame, spatial \lewm{} maps green/orange to olive
    while preserving layout; \pldm{} and \lpwm{} retain more separation.
    \dinowm{} is clearest but uses an unmatched minibatch ($\dagger$).  Bottom:
    spatial \lewm{} (IDM .125, chosen to make the counterexample legible; the
    full six-weight grid is in the appendix) separates three TR3 colors, whereas
    swapping orange for blue in TR2 yields teal.  All-marker audits and probes provide descriptive
    corroboration beyond these qualitative panels.}
    \label{fig:teaser}
\end{figure*}

\section{Introduction}

Color is a nuisance only when it is irrelevant to dynamics.  In Tunnel Rooms
Color, portals share geometry but their colors determine different
destinations; predicting a portal transition therefore requires information
about the color--location assignment.  Yet across the tested random-layout
two-room settings, spatial LeWorldModel (\lewm{}; \citealp{maes2026leworldmodel})
reconstructs green and orange as the same olive marker while accurately
localizing both portals (\cref{fig:teaser}).

JEPAs predict target representations rather than reconstruct pixels
\citep{lecun2022path,assran2023ijepa,bardes2024vjepa}.  Here the markers remain
visible, but their consequence occurs only at a crossing.  This creates a
possible availability conflict: average prediction can favor geometry and
local motion over color identity.  Related toy JEPAs can retain a persistent
distractor while losing a moving dot \citep{sobal2022slow}; our persistent cue
instead determines the next state.  We test the failure but do not identify its
optimization mechanism.

We introduce \emph{Tunnel Rooms Color}: an agent moves from a hub through two
or three colored portals, each paired with a same-color destination.  Geometry
and local motion are predictable without hue; teleportation identity is not.
We call a state operationally ``color blind'' when specified readouts cannot
recover hue/binding while geometry remains accessible---not when
$I(Z;C)=0$.

Across native CLS and spatial 8$\times$8 representations and six
inverse-dynamics-modeling (IDM) weights, TR2 LeWM decoders reconstruct the cue
toward a mixed color.  Two spatial green/blue controls at IDM 0 and .625 repeat
the result.  A stagewise audit of all 12 TR2 checkpoints finds hue in native
patches but weak accessibility in every projected state; exact marker swaps
and nonlinear readouts corroborate one prespecified spatial checkpoint.

Three-seed readout controls sharpen the boundary.  LeWM is weak under native
CLS, flattened-patch MLP, and spatial 8$\times$8 global states.  PLDM
\citep{sobal2025latentplanning} improves only with spatial 8$\times$8, whereas
a Product-Laplace LpWM adaptation of Rectified LpJEPA
\citep{kuang2026rectified} is already substantially better with native CLS.
Thus neither comparator is uniformly color-sensitive, but each has a strong
TR2 global readout absent from LeWM.  At selected spatial checkpoints, PLDM
retains visible hue but not destination binding, whereas LpWM retains both.
Native all-patch \dinowm{}
\citep{zhou2025dinowm} reconstructs color, and spatial LeWM recovers color in
TR3.  These controls establish observation-level solvability and important
boundary conditions; objective and tokenization differences preclude a pure
architecture ranking.

\section{Tunnel Rooms Color Separates Geometry from Identity}

\paragraph{Environment.}
Each 224$\times$224 RGB observation shows a central hub, an agent, and two
(\trtwo{}) or three (\trthree{}) destination rooms (leaves).  Colored markers
pair every hub portal with the same-color leaf reached after entry.  Each
episode specifies a target leaf; the correct hub portal shares its color,
while randomized layouts prevent solving identity from a fixed position.
Portal crossings account for 1.90\%/2.12\%
of adjacent transitions in the audited \trtwo{}/\trthree{} corpora, making the
color-conditioned consequence less frequent even though the marker remains
visible.  Under the exact TR2 four-observation, frameskip-five loader, 10.22\%
of legal windows contain a teleport in the prediction target interval and
26.62\% contain one anywhere in context or target.

\paragraph{Models and factorial.}
We train \lewm{} with IDM weights
$\{0,.005,.025,.125,.625,3.125\}$, native CLS or spatial 8$\times$8
aggregation, and \trtwo{} or \trthree{}, giving 24 runs.  We compare one
spatial \pldm{}, spatial \lpwm{}, and native all-patch \dinowm{} per dataset.
All use seed 3072 and ten epochs.  The first three use trainable ViT-tiny
encoders; their spatial variants use a shared 8$\times$8 aggregator, whereas
the other \lewm{} variants use native CLS.  \dinowm{} uses frozen pretrained
DINOv2 patches, a patch decoder, and more updates
\citep{oquab2023dinov2,zhou2025dinowm}.
Equal epochs yield about 32,550 TR2 versus 20,820 TR3 updates, so room count is
not an isolated intervention.

Two follow-ups test global readout and temporal smoothness.  A three-seed TR2
grid compares native CLS, flattened-patch MLP, and spatial 8$\times$8 states
within LeWM, PLDM, and LpWM.  Each family retains its labeled objective recipe,
so within-family readout comparisons are controlled and family comparisons are
boundary conditions.  A one-seed LeWM sweep adds PLDM-style adjacent-state MSE
at weights $\{0,.005,.025,.125,.2,.625,3.125\}$; .2 matches PLDM.

For observation $x_t$, let
\begin{equation}
 z_t=g_\phi(f_\theta(x_t)), \qquad
 \hat z_{t+1}=p_\psi(z_{\leq t},a_{\leq t}).
 \label{eq:jepa}
\end{equation}
The audited LeWM+IDM objective has the schematic form
\begin{equation}
 \mathcal L_{\rm audit}=\mathbb E\|\hat z_{t+1}-z_{t+1}\|_2^2
 +\lambda\,\mathrm{SIGReg}(Z)+\beta\mathcal L_{\rm IDM}.
 \label{eq:lewm-objective}
\end{equation}
The original LeWM objective is recovered when $\beta=0$
\citep{balestriero2025lejepa,maes2026leworldmodel}.  Standardized ridge probes
fit on 640 examples and test on 256 probe-disjoint examples.  Semantic targets
locate hubs by color; positional controls locate fixed geometric slots.  Their
bias-only references are 30.02/19.27 pixels (TR2) and 30.25/7.72 (TR3).
Semantic MAE therefore combines identity and localization.

A separate 525/75/150 episode split localizes fixed-slot hue and
target--destination binding across native patches, nonlearned pooling, learned
aggregation, and projection.  Each stage selects its ridge penalty on
validation; 50\% is balanced chance.  One pre-frozen LeWM checkpoint supplies
the deeper nonlinear and counterfactual controls.  Full probe and split
contracts appear in the appendix.

\paragraph{Why averaging is diagnostic.}
For an MSE decoder $d_\omega$, the population optimum is
\begin{equation}
 d^*(z)=\mathbb E[X\mid Z=z].
 \label{eq:conditional-mean}
\end{equation}
If $Z$ does not distinguish cue $C$ at the marker pixels, then
$d^*(z)\approx\sum_c p(c\mid z)x_c$: green/orange becomes olive and
green/blue becomes teal.  This signature is \emph{consistent with} cue
collapse, not proof of $I(Z;C)=0$; decoder capacity is an alternative.
Conditional-mean blur under MSE is documented in pixel prediction
\citep{mathieu2016video}.  Representations are detached, so decoder training
cannot make the world model retain color.

From each final video we detect marker support in the original, exclude agent
overlap, and sample the matched decode.  We report CIE L$^*$a$^*$b$^*$
nearest-centroid accuracy and decoded/original class-centroid separation.
Support-sensitivity checks preserve the ordering; clips are descriptors, not
independent replicates.

\section{LeWM Preserves Geometry but Misses Color in TR2}

\begin{figure*}[t]
    \centering
    \includegraphics[width=\textwidth]{figures/color_sweep_summary.pdf}
    \caption{\textbf{LeWM's TR2 color failure persists across global readouts
    and temporal alignment.}  (a,b) Three-seed mean $\pm$ SD: all LeWM color
    readouts remain weak despite accurate spatial geometry; PLDM spatial and
    LpWM CLS are stronger.  (c,d) One-seed temporal sweeps: no weight rescues
    TR2; spatial TR3 improves while CLS remains near null.  Dotted lines are
    target-specific bias-only references.}
    \label{fig:sweep-summary}
\end{figure*}

\paragraph{Geometry improves; color-conditioned localization does not.}
In the three-seed no-IDM replication, LeWM's current semantic hub MAE is
$30.69\pm.47$ for native CLS, $28.35\pm.74$ for MLP, and $25.97\pm.20$
pixels for spatial 8$\times$8.  Position-only MAE falls from $19.48\pm.29$ to
$13.41\pm.61$ and $6.41\pm.34$ pixels over the same readouts
(\cref{fig:sweep-summary}).  The spatial representation therefore removes
66.7\% of its positional null error but only 13.5\% of its semantic null error.
Across the earlier six-weight spatial IDM sweep, semantic MAE likewise stays
24.84--25.79 pixels while position-only MAE reaches 4.50 pixels.  IDM improves
geometry without making color identity reliably accessible.

\paragraph{Temporal alignment does not rescue TR2.}
Across weights 0--3.125, current TR2 semantic hub MAE spans only
25.74--25.84 pixels for spatial 8$\times$8 and 29.89--30.96 for CLS;
one-step predictions closely overlap.  The PLDM-matched .2 coefficient is no
exception.  In TR3, however, spatial MAE improves from 9.00 to 7.58 pixels as
the weight increases, while CLS stays near 29 pixels.  Temporal smoothness is
therefore neither sufficient to create a useful TR2 global color state nor a
universal suppressor of color.  Because this mechanism sweep uses one seed,
the TR3 dose response is descriptive.

\paragraph{Fixed-slot visible hue remains weak across the factorial.}
Balanced accuracy spans 45.0--50.0\% for CLS and 49.2--56.3\% for spatial
projectors; target--destination binding spans 46.2--51.3\% across all 12.
Every corresponding 95\% Wilson interval for raw accuracy includes 50\%, and
neither endpoint varies monotonically with IDM.  All cells share one
world-model seed and test split: this is a hyperparameter-robustness check, not
training-seed replication.

\paragraph{Hue accessibility declines after the patch representation.}
The stagewise audit localizes the TR2 failure across all 12 checkpoints
(\cref{tab:tr2-stagewise}).  Native ViT patches classify fixed-slot hue at
96.0--100\% balanced accuracy, and nonlearned 8$\times$8 pooling reaches 100\%
throughout.  The actual CLS state is already weak, while learned spatial
aggregation retains 66.0--92.7\%; projection reduces every checkpoint to the
near-chance range reported above.  The native-to-projector decline is positive
in 12/12 configurations and averages 50.1 percentage points.  Thus backbone patch tokens
retain visible hue; accessibility declines along the learned global path and
is weak in every final projector.  On the prespecified spatial checkpoint, the
learned aggregator reaches 66.0\%.  Inside its projector, hue stays at
65.8--68.7\% through GELU but falls to 51.6\% after the final
2048$\rightarrow$192 Linear.  A fixed 192-D map retains 71.1\%, while selected
unmatched PLDM stays at 100\% throughout.

\begin{table}[h]
\centering
\footnotesize
\setlength{\tabcolsep}{2.0pt}
\begin{tabular}{@{}llrrrr@{}}
\toprule
Model & Stage & Hue & Bind. & Portal & Dec. sep. \\
\midrule
\lewm{} & native patches & 96.0 & 48.0 & 58.2 & -- \\
\lewm{} & projector & 51.6 & 50.6 & 52.1 & 4.85 \\
\pldm{} & projector & 100 & 46.4 & 52.6 & 49.5 \\
\lpwm{} & projector & 100 & 95.2 & 86.9 & 37.8 \\
\dinowm{} & native patches & 100 & 60.2 & 54.0 & 78.8 \\
\bottomrule
\end{tabular}
\caption{\textbf{Where TR2 color survives in selected checkpoints.}
Hue, binding, and portal columns are balanced accuracy (\%); chance is 50\%.
Decoder separation is the percentage of original green--orange centroid
separation retained on 300 prespecified frames.  Comparator checkpoints are
selected and unmatched; DINO-WM has no global projector.  The complete
12-checkpoint LeWM hue audit appears in \cref{fig:stagewise-factorial}.}
\label{tab:tr2-stagewise}
\end{table}

The final LeWM state does not exceed measured geometry baselines for binding
(50.6\%) or portal identity (52.1\%).  At these unmatched checkpoints, only
LeWM's global projector is near chance on hue; PLDM retains hue but not binding, whereas LpWM
retains hue/binding/portal identity (100/95.2/86.9\%).  A
validation-selected one-hidden-layer MLP likewise leaves LeWM hue/binding at
52.4/47.0\%, but LpWM at 100/96.9\% (five probe seeds).  This excludes only
that MLP family; exact intervention results are in \cref{sec:counterfactual}.
At this audited checkpoint, near-unit variance and accurate geometry reject
total scalar collapse.  Because the current projected state already has a weak
hue readout, the predictor is not the earliest audited forward stage with weak
accessibility; predictor-loss gradients during training remain unresolved.

\paragraph{The decoder moves toward the palette mean.}
For spatial \trtwo{} at IDM 0--.625, final-video color accuracy is 50\% and
only 4.0--7.5\% of original separation remains.  On 300 prespecified frames,
LeWM retains 4.85\% separation (95\% CI [3.57\%, 6.17\%]), versus 49.5/37.8/
78.8\% for PLDM/LpWM/DINO-WM; color accuracies are 54.2/57.3/70.4/100\%.
Thus PLDM retains a color axis despite modest class accuracy, while LpWM has
greater class recoverability than LeWM.  These controls are selected and
unmatched.  LeWM has lower pixel MSE than PLDM/LpWM (.0996 versus
.1190/.1305), so average reconstruction error hides color failure.  Green/blue
controls likewise render teal and retain only 2.9/2.1\% separation (appendix).

\section{The Failure Is Conditional, Not Universal}

\paragraph{Comparator success depends on the global readout.}
LeWM is the only family in the three-seed grid without a strong TR2 readout:
its best semantic mean is $25.97\pm.20$ pixels.  PLDM reaches
$18.05\pm.41$ with spatial 8$\times$8, although its CLS is near null and its
MLP is seed-unstable.  LpWM reaches $15.65\pm1.42$ with native CLS, although
its MLP is weak and its spatial results split sharply across seeds
(11.95/28.36/28.66 pixels).  The claim is therefore not that every PLDM or
LpWM state preserves color; rather, both families expose a substantially
better global state than any tested LeWM readout.  Native all-patch \dinowm{}
also decodes color at 100\% accuracy with 81.9/76.5\% TR2/TR3 separation.

Selected stagewise controls agree: PLDM retains visible hue but not binding,
whereas LpWM retains both (\cref{tab:tr2-stagewise}).  A VISReg intervention
rescues one seed but not two others or binding (\cref{tab:visreg-swap}); the
objective's causal role remains unresolved.

\paragraph{TR3 is a boundary condition.}
Spatial \lewm{} reaches 7.58-pixel semantic MAE and 79--97\% decoder accuracy
on \trthree{}; CLS remains near null.  This refutes universal blindness but,
because exposure differs, does not show that more colors cause recovery.

\clearpage
\bibliography{references}
\bibliographystyle{icml2026}

\clearpage
\appendix

\section{Three-Seed Global-Readout Replication}

\Cref{tab:aggregation-replication} reports the exact TR2-Color probe summary
behind \cref{fig:sweep-summary}a--b.  ``Position'' is a geometric control
measured on the same color dataset, not a separate non-color benchmark.  The
per-seed rows and color-specific green/orange errors are frozen in
\texttt{data/tr2\_aggregation\_color\_probe.csv}.

\begin{table*}[t]
\centering
\small
\setlength{\tabcolsep}{4.2pt}
\begin{tabular}{@{}llrrrr@{}}
\toprule
& & \multicolumn{2}{c}{Color-conditioned hub MAE} &
\multicolumn{2}{c}{Position-only hub MAE} \\
\cmidrule(lr){3-4}\cmidrule(lr){5-6}
Model & Global readout & Current & Predicted & Current & Predicted \\
\midrule
\lewm{} & CLS & $30.69\pm.47$ & $30.69\pm.50$ & $19.48\pm.29$ & $19.49\pm.31$ \\
\lewm{} & MLP & $28.35\pm.74$ & $28.35\pm.62$ & $13.41\pm.61$ & $13.43\pm.62$ \\
\lewm{} & spatial 8$\times$8 & $25.97\pm.20$ & $26.02\pm.22$ & $6.41\pm.34$ & $6.44\pm.32$ \\
\midrule
\pldm{} & CLS & $31.05\pm.32$ & $30.67\pm.33$ & $20.09\pm.04$ & $19.78\pm.11$ \\
\pldm{} & MLP & $28.57\pm5.02$ & $28.75\pm5.08$ & $20.03\pm1.17$ & $19.86\pm1.24$ \\
\pldm{} & spatial 8$\times$8 & $18.05\pm.41$ & $17.17\pm.29$ & $16.54\pm.50$ & $16.23\pm.32$ \\
\midrule
\lpwm{} & CLS & $15.65\pm1.42$ & $15.74\pm1.62$ & $11.26\pm.51$ & $11.27\pm.45$ \\
\lpwm{} & MLP & $29.79\pm.69$ & $28.99\pm1.18$ & $13.47\pm3.50$ & $12.92\pm3.48$ \\
\lpwm{} & spatial 8$\times$8 & $22.99\pm9.56$ & $23.09\pm10.24$ & $12.42\pm.77$ & $12.35\pm1.16$ \\
\bottomrule
\end{tabular}
\caption{\textbf{TR2-Color global-readout replication.} Mean $\pm$ one sample
standard deviation over seeds 3072/4072/5072; lower is better.  Each family
retains its labeled recipe: LeWM uses SIGReg without IDM, PLDM uses VCReg with
temporal alignment .2, and LpWM uses Product-Laplace InfoReg with IDM 3.125.
The comparison isolates readout within each family, not objective across
families.}
\label{tab:aggregation-replication}
\end{table*}

\section{Four-Color Boundary Experiment}

The spatial \trthree{} result suggests that greater cue diversity may make
color harder for LeWM to ignore, but it does not isolate color cardinality:
TR2 and TR3 differ in geometry, trajectory distribution, transition exposure,
and optimization steps.  A four-color result has the same ambiguity, so we
treat \trfour{} as a prespecified boundary experiment rather than a causal
cardinality ablation.

The randomized environment has four right-side leaf rooms and four unique
portal identities (green, orange, blue, and magenta).  Portal positions,
pairings, cue assignments, starts, and targets are sampled per episode.  The
comparison trains spatial 8$\times$8 LeWM, PLDM, and LpWM for ten epochs over
three seeds, with each family retaining the labeled recipe used in
\cref{tab:aggregation-replication}.  The primary endpoints are mean
color-conditioned hub MAE over all four identities relative to a target-mean
reference, position-only hub MAE, and decoded centroid-separation retention.
Fixed-slot hue and destination binding are complementary; ranked planning is
secondary because visible hue without correct binding is insufficient for
control.

The 10,000-episode corpus completed with 638,774 frames and an approximately
uniform geometry-defined rank histogram (2,438/2,553/2,509/2,500). The first
dependency-released training submission exited before model initialization:
its source-commit guard correctly detected that the shared checkout had
advanced. A second submission then exposed a missing Hydra schema entry for
the new four-leaf protocol, also before model initialization. We therefore
report no TR4 model result from either operational attempt. The same
prespecified nine-run matrix was corrected without changing its scientific
design and submitted as revision 4 from an immutable detached worktree; its
outcomes remain pending.
\experimentneeded{After all revision-4 checkpoints finish, freeze the decoder,
location-probe, hue, and binding audits before adding any TR4 outcome.}
If spatial LeWM remains weak while PLDM/LpWM retain four-way identity, the TR2
bottleneck extends despite greater cue diversity.  If LeWM improves, the result
locates another boundary like TR3 but cannot attribute it to color count alone.
Repeating two colors across four portals is not a matched control: it makes
color intrinsically insufficient to identify a unique destination.  A causal
follow-up must preserve four unique identities while separately controlling
palette separation, geometry, transition exposure, and optimization updates.

\section{Complete LeWM Factorial}

The full numeric table is distributed as
\texttt{data/lewm\_color\_identity\_factorial.csv}.  It contains every stable
anonymous experiment ID, aggregation, IDM weight, room count, current/predicted semantic MAE,
position-only MAE, their raw difference, and seed.  The raw difference is
descriptive, not a pure hue statistic; target-specific bias-only references
are in \texttt{data/probe\_target\_nulls.csv}.  No missing value is encoded as zero.
All 24 runs completed ten epochs and their final diagnostics.

\begin{figure*}[t]
    \centering
    \includegraphics[width=\textwidth]{figures/lewm_identity_factorial.pdf}
    \caption{\textbf{Complete original LeWM IDM factorial.} Mean
    color-conditioned and position-only hub MAE for current and one-step
    predicted states.  Spatial TR2 improves geometry but not color binding;
    spatial TR3 is the counterexample.  All cells share seed 3072.}
    \label{fig:idm-factorial}
\end{figure*}

\begin{table}[h]
\centering
\small
\setlength{\tabcolsep}{3.2pt}
\begin{tabular}{@{}rrrrr@{}}
\toprule
IDM & CLS hue & CLS bind. & 8$\times$8 hue & 8$\times$8 bind. \\
\midrule
0     & 47.7 & 48.0 & 50.0 & 50.0 \\
.005  & 50.0 & 48.0 & 50.0 & 47.2 \\
.025  & 45.0 & 46.2 & 53.6 & 49.9 \\
.125  & 47.9 & 48.9 & 56.3 & 48.3 \\
.625  & 47.9 & 48.9 & 49.2 & 47.0 \\
3.125 & 45.3 & 51.3 & 51.6 & 50.6 \\
\bottomrule
\end{tabular}
\caption{LeWM TR2 projector balanced accuracy (\%) on fixed-slot hue and
target--destination color binding.  All cells share training seed 3072 and the
same frozen 525/75/150 episode split; 50\% is balanced chance and no
population-level training-seed uncertainty is estimated.}
\label{tab:projector-factorial}
\end{table}

\begin{figure*}[t]
    \centering
    \includegraphics[width=\textwidth]{figures/lewm_tr2_stagewise_hue_factorial.pdf}
    \caption{\textbf{The stagewise hue decline holds across every audited TR2
    LeWM checkpoint.}  Each point is balanced accuracy on the same 150
    readout-test episodes; lines connect hyperparameters for readability, not
    training trajectories.  Native patches and deterministic 8$\times$8
    pooling retain hue; the latter is a readout control and is not part of the
    CLS computation path.  The actual CLS state is weak, whereas learned spatial
    aggregation retains varying hue; every projector is near chance.  All
    checkpoints share one world-model seed, so this is a hyperparameter audit,
    not training-seed replication.}
    \label{fig:stagewise-factorial}
\end{figure*}

\section{Tunnel Rooms Color Specification}

TR2-Color and TR3-Color each contain 10,000 successful expert episodes with
651,357 and 482,051 stored RGB observations.  The agent receives a continuous
two-dimensional motion action, repeated for five simulator steps per stored
transition.  Each reset randomizes the agent, goal, portal geometry, and a
permutation assigning two or three cue identities to the hub/leaf portal
pairs.  Entering a hub portal teleports the agent to its paired leaf; the goal
specifies a leaf destination, so success requires the same-color assignment
rather than the nearest portal.  Median episode lengths are 56/45 frames.
Using a within-episode position jump above ten pixels as the prespecified
crossing detector, portal transitions are 12,180/641,357 (1.90\%) in TR2 and
10,000/472,051 (2.12\%) in TR3.  The world-model 90/10 split is over legal
short windows, not entire episodes; the stagewise probe split is
episode-disjoint, but it draws from this source corpus.  Exact counts,
definitions, and dataset hashes are frozen in
\texttt{data/benchmark\_audit.csv}.  The maintained loader yields target
intervals containing a crossing in 10.22\% of all legal windows and 10.21\% of
the exact 90\% training split; counts and split seed are frozen in
\path{data/training_window_transition_audit.csv}.

\section{Decoder Color-Fidelity Audit}

\paragraph{Protocol.}
For each of 34 final-epoch training-media pairs (24 \lewm{}, six selected
comparators, two native DINO CLS controls, and two blue/green controls), we
parse six 224$\times$224 clips over four
frames.  Saturated 18$\times$18 non-red connected components in the original
identify portal support; red connected components identify the agent, and any
overlap is excluded.  The same support is sampled from the decoded frame.
Run-local original CIE L$^*$a$^*$b$^*$ centroids define the nearest-color rule
and the unweighted palette mean.  Worst matched-pair separation retention is
\begin{equation}
 R_{\mathrm{sep}}=
 \min_{c\ne c'}
 \frac{\|\hat\mu_c-\hat\mu_{c'}\|_2}
      {\|\mu_c-\mu_{c'}\|_2}.
 \end{equation}
This is a reconstruction-at-known-support diagnostic, not a test of mutual
information.  Exact media hashes, instance counts, exclusions, and run-level
statistics are preserved with the paper.

\paragraph{Prespecified diagnostic episodes.}
We also apply the same metric to two marker-visible frames from each of 150
distinct episodes in the stagewise test split; selections were fixed before
decoder outputs were inspected.  The 300 frames
yield 1,127 cue instances after excluding 72 agent overlaps and one failed
detection.  Checkpoints and decoder sidecars are strict-loaded and decoder
tensors are checked bitwise.  A 10,000-replicate episode bootstrap gives
separation retention of 4.85\% [95\% CI 3.57\%, 6.17\%] for LeWM,
49.5\% [48.4\%, 50.5\%] for PLDM, 37.8\% [36.4\%, 39.1\%] for LpWM, and
78.8\% [78.1\%, 79.5\%] for DINO-WM.
LeWM's reconstructed green/orange centroids are nearly identical despite its
lower normalized-pixel MSE than PLDM/LpWM.  These prespecified episodes come
from the world-model source corpus and are therefore not a world-model
holdout; logged media lacks source IDs, so disjointness from its six clips
cannot be proved.  An anonymized frame manifest and all cue-level sufficient
statistics are distributed for independent recomputation.

For the frozen LeWM checkpoint, all four trained decoder snapshots logged over
training retain only 4.0--9.4\% of palette separation; the random-decoder
snapshot is excluded.  The snapshots use different minibatches and W\&B media
indices are not epochs, so this rules out only an obvious late collapse from a
previously high-fidelity decoder.  It is not an epoch-wise representation
probe; portable intermediate checkpoints were not saved.

\section{Stagewise Hue and Binding Readouts}

The primary \trtwo{} audit draws 750 episodes from the world-model training
corpus.  Its cue tasks use exactly one example per episode: the goal image and
target position stored at the episode offset.  Two separately sampled frames
per episode are used only by an auxiliary agent-position probe.  Fit,
validation, and test episodes are disjoint (525/75/150), and every target
selects its own ridge penalty on validation before a final refit on the 600
fit-plus-validation episodes.  For ridge readouts, classification penalties maximize raw validation accuracy;
balanced accuracy is used only for test reporting.  Features are centered per
coordinate and divided by one training-set global RMS scalar; coordinate
targets are divided by 224 and class targets are one-hot vectors.  Targets are
centered on the ridge-training split and the mean is restored after prediction.
No examples or classes are reweighted.

Visible hue is the identity of the upper leaf marker; both hues occur in every
image, so this is a spatial identity task rather than color-presence detection.
Destination binding is the hue of the leaf nearest the stored target.  The
portal readout regresses the coordinate of the hub with that hue, then assigns
the prediction to the nearest actual hub and scores its color.  A
room-majority destination-color rule and a global-mean hub coordinate obtain
54.8\% and 54.5\% balanced accuracy on the same test split.  These are measured
geometry baselines, not statistical null distributions; theoretical balanced
chance remains 50\%.

As a nonlinear check, one-hidden-layer ReLU MLP candidates use fit-split
per-feature standardization and are fit on the 525 training episodes.  Width $\{32,64,128\}$, weight decay
$\{0,10^{-4},10^{-3}\}$, and an epoch up to 300 are ranked by validation balanced
accuracy, then, in order, lower class-weighted validation cross-entropy,
smaller width, smaller weight decay, and earlier epoch.  The selected recipe is refit on all
600 fit-plus-validation episodes and tested once on 150 for each of five
optimizer seeds.
LeWM hue/binding remains 52.4$\pm$3.0\%/47.0$\pm$3.7\% balanced accuracy; the
20,000-replicate episode-bootstrap intervals are 47.6--57.3\% and
43.0--51.0\%.  PLDM retains hue but not binding, while LpWM retains both.
Probe seeds quantify optimization variability, not world-model seed
uncertainty.

The projector-substage audit uses the same split and independently selects a
ridge penalty at every stage:

\begin{table}[h]
\centering
\small
\setlength{\tabcolsep}{3.5pt}
\begin{tabular}{@{}lrrrrr@{}}
\toprule
Model & Input & Linear$_1$ & BN & GELU & Linear$_2$ \\
\midrule
\lewm{} & 66.0 & 65.8 & 68.7 & 67.0 & 51.6 \\
\pldm{}$^\dagger$ & 100 & 100 & 100 & 100 & 100 \\
\bottomrule
\end{tabular}
\caption{Fixed-slot hue balanced accuracy (\%) within one 192--2048--192
projector.  Linear$_2$ exactly reproduces each prior full-projector result;
$\dagger$PLDM is a selected unmatched control.}
\label{tab:projector-substage}
\end{table}

Exact all-marker and hub-only green/orange swaps alter 972,000 and 486,000
exact-RGB pixels of the filled 18$\times$18 markers over the full audit,
respectively, with zero changes elsewhere.  The all-marker swap flips visible
and destination-color labels and reverses color-to-hub indexing while leaving
the correct physical hub coordinate fixed.  The hub-only swap leaves visible
and destination labels fixed, reverses hub-color indexing, and moves the
correct-hub coordinate.  The complete stagewise values,
selected penalties, Wilson intervals, counterfactual readouts, checkpoint
hashes, and split hash are frozen in the accompanying \texttt{tr2\_stagewise}
tables.  Because the world model trained on the source corpus, this establishes
readout generalization across episodes, not world-model generalization to a
new corpus.

\subsection{Counterfactual binding sensitivity}
\label{sec:counterfactual}

We sample one three-frame context from each of 128 distinct TR2 episodes and
compare three exact interventions: exchange the two marker RGB values, shift
both markers by a common RGB offset, or recolor the agent.  The first two alter
the same 3,888 pixels and differ in input relative $\ell_2$ magnitude by less
than .5\%.  For representation stage $s$, define
\begin{equation}
\begin{aligned}
 \delta_s(v)&=\|h_s(x^v)-h_s(x)\|_2/\|h_s(x)\|_2,\\
 Q_s&=\delta_s({\rm swap})/\delta_s({\rm shift}),\\
 H_s&=Q_s/Q_{\rm native}.
\end{aligned}
\end{equation}
$Q_s$ measures response to exchanging color--location binding relative to a
matched generic color perturbation; $H_s$ measures its retention from native
patches.  Values are paired episode means with 20,000-replicate bootstrap
intervals.

\begin{table}[h]
\centering
\small
\setlength{\tabcolsep}{2.5pt}
\begin{tabular}{@{}lrrrr@{}}
\toprule
Model & $Q_{\rm native}$ & $Q_{\rm agg}$ & $Q_{\rm proj}$ & $Q_{\rm pred}$ \\
\midrule
\lewm{}, IDM 3.125 & .855 & .341 & .287 & .280 \\
\lewm{}, IDM .625 & .605 & .147 & .134 & .133 \\
\pldm{} & 1.956 & 2.898 & 3.197 & 3.765 \\
\lpwm{} & 1.235 & .826 & 1.769 & 2.340 \\
\dinowm{} & .528 & .528 & -- & .574 \\
\bottomrule
\end{tabular}
\caption{TR2 counterfactual response relative to the matched cue-shift control.
For DINO-WM, aggregation is the unchanged native patch grid. Predictor/native
retention is .327/.219 for the two LeWM checkpoints, 1.908 for PLDM, 1.896 for
LpWM, and 1.085 for DINO-WM. Full 95\% intervals and the agent control are in
\texttt{data/tr2\_counterfactual\_sensitivity.csv}.}
\label{tab:counterfactual}
\end{table}

The accompanying raw-response table reports numerator and denominator
separately: $Q_s$ can change because either response changes, so it is a
relative selectivity statistic rather than proof that absolute swap
sensitivity is erased.  Agent recoloring is retained only as an unmatched
descriptive calibration and is not interpreted.  One
seed, a source-corpus audit, and an out-of-distribution common-shift control
limit causal generalization; absolute responses are not compared across
latent scales.

For context, the same stagewise visible-hue readout protocol on \trthree{}
provides a boundary condition:

\begin{table}[h]
\centering
\small
\setlength{\tabcolsep}{2.5pt}
\begin{tabular}{@{}lrrrrr@{}}
\toprule
Model & Native & Pool & Rand. & Agg. & Proj. \\
\midrule
\pldm{} & 100.0 & 100.0 & 100.0 & 100.0 & 100.0 \\
\lewm{} & 99.3 & 100.0 & 74.7 & 89.3 & 80.7 \\
\lpwm{} & 100.0 & 100.0 & 100.0 & 100.0 & 100.0 \\
\dinowm{} & 99.3 & -- & -- & -- & -- \\
\bottomrule
\end{tabular}
\caption{Visible-cue color accuracy (\%) across \trthree{}
representation
stages.  \lewm{} is the 8$\times$8, IDM-3.125 run; \lpwm{} is the later
8$\times$8, IDM-.625, temporal-Jaccard-.1 run.  These unmatched settings make
the table a localization clue, not an architecture ranking; each row uses 150
test episodes.}
\label{tab:stagewise-tr3}
\end{table}

The TR3 LeWM projector reaches 80.7\% raw accuracy, versus 53.3\% TR2 raw
accuracy under the matched readout protocol.  Native-token hue is accessible
in both tasks; the learned
global path's retention is conditional on the environment.  Different task
statistics, 50\% more marker-support pixels per frame, and unequal update
counts preclude attributing the contrast to room count alone.

\section{Alternate Palette and Fixed-Layout Control}

\paragraph{Fixed layouts can masquerade as color understanding.}
With corrected episode-level cue assignments, the same LeWM family obtains
1.81-pixel post-projector color-hub MAE and $R^2=.994$ in a three-room
fixed-wall pilot, versus 17.74 pixels when cue geometry varies.  Because a
fixed layout makes the semantic target recoverable from position, this result
does not test hue accessibility.  It instead demonstrates why the primary audit
randomizes marker geometry and pairs semantic readouts with positional controls.

The green/blue controls use fixed anonymous experiment records at IDM 0 and .625 and the same seed
3072.  The source-dataset SHA256 is recorded in the paper data with prefix
\texttt{6e2823ad}.
Their full probe values appear in \texttt{data/blue\_green\_control.csv}; visual
examples appear in \cref{fig:teaser}.

\section{LpWM Sparsity and Temporal-Jaccard Controls}

We backfill the identical 640/256 location-probe protocol on an internally matched
single-seed LpWM matrix with Product-Laplace
$\mu\in\{0,-0.5,-1,-1.5\}$, temporal-Jaccard weight 0 or .1, fixed $p=1$,
IDM .625, and spatial 8$\times$8 features.  Across all eight checkpoints,
current semantic hub MAE is 8.25--11.10 pixels and positional hub MAE is
8.03--11.17 pixels (\cref{tab:lpwm-tj}).  Thus LpWM does not reproduce LeWM's
large semantic-versus-position gap in this controlled family.  Temporal
Jaccard shows no directionally consistent benefit: its same-seed, same-split contrast
changes sign with $\mu$, worsening both hub targets at $\mu=-1$, improving both
at $\mu=-1.5$, and producing mixed effects otherwise.  All cells share one
training seed, so these are conditional checkpoints rather than a
population-level factorial result.

\begin{table}[h]
\centering
\scriptsize
\setlength{\tabcolsep}{3.1pt}
\begin{tabular}{@{}rrrrr@{}}
\toprule
$\mu$ & TJ & Sem. hub & Pos. hub & Pred. sem. \\
\midrule
0 & 0 & 8.93 & 8.03 & 8.87 \\
0 & .1 & 8.70 & 8.13 & 8.92 \\
-.5 & 0 & 8.96 & 9.17 & 8.98 \\
-.5 & .1 & 8.50 & 10.24 & 8.68 \\
-1 & 0 & 8.25 & 8.40 & 8.32 \\
-1 & .1 & 9.14 & 9.47 & 9.40 \\
-1.5 & 0 & 11.10 & 11.17 & 11.49 \\
-1.5 & .1 & 10.48 & 10.85 & 10.67 \\
\bottomrule
\end{tabular}
\caption{\textbf{Internally controlled single-seed LpWM TR2-Color grid.} Current-feature semantic and
position-only hub MAE and predicted-feature semantic hub MAE (pixels).  Lower
is better.  All rows use seed 3072, the same 256-episode probe-test split, and
fixed ridge $\alpha=10$; that split is held out from probe fitting, not
world-model training.  TJ pairs are distinct training runs, so no across-seed
uncertainty is estimated.}
\label{tab:lpwm-tj}
\end{table}

\section{Exploratory Shape-Conditioned Localization Caveat}

We replaced color with identically colored triangle, diamond, and star markers
while preserving the underlying trajectories, actions, positions, rewards, and
episode boundaries.  From the same frozen current projected world-model states
and one-step predicted states, ridge probes predict either the location of a
named shape or a positional slot.  Across one seed, named-shape MAE is
27.02--31.18 pixels for current states and
26.84--30.76 for predictor features, near intercept-only baselines of 31.16
(TR2) and 30.20 (TR3).  Current and predicted features support substantially
better positional than shape-conditioned localization
(\cref{tab:shape-control}).  Because PLDM and the historical learned-token
DINO-WM control show the same pattern, and no completed LpWM Shape result exists, this pilot
is a broader cue-binding caveat rather than evidence for extending the main
LeWM-specific color claim.

Final-epoch decoder clips agree qualitatively but are not an independent
identity metric. Across all available clips, LeWM and PLDM reconstruct portal
locations and often the agent, yet triangle, diamond, and star markers soften
into similarly colored diffuse blobs; learned-token DINO-WM is blurrier and
often suppresses the agent as well. No existing family visibly reconstructs
the three identities consistently enough to support a LeWM-specific decoder
failure. This assessment uses every final-epoch clip rather than one selected
frame, but pixel-MSE decoder capacity remains a confound.

\begin{table}[h]
\centering
\scriptsize
\setlength{\tabcolsep}{3.0pt}
\begin{tabular}{@{}llrrrr@{}}
\toprule
Data & Model & \multicolumn{2}{c}{Identity MAE} & \multicolumn{2}{c}{Position MAE} \\
& & Current & Pred. & Current & Pred. \\
\midrule
TR2 & DINO--MLP & 27.02 & 26.84 & 7.47 & 7.36 \\
TR2 & PLDM & 29.17 & 28.37 & 6.80 & 6.69 \\
TR2 & LeWM & 28.05 & 28.10 & \textbf{4.57} & \textbf{4.41} \\
TR3 & DINO--MLP & 27.93 & 27.87 & 6.94 & 6.85 \\
TR3 & PLDM & 31.18 & 30.76 & 3.18 & 3.20 \\
TR3 & LeWM & 29.99 & 29.80 & \textbf{2.66} & \textbf{2.56} \\
\bottomrule
\end{tabular}
\caption{\textbf{Exploratory shape-conditioned localization control (pixels).} Lower is better.
Identity probes locate a named shape; position probes locate a geometric slot.
All rows use one world-model seed.  DINO--MLP is a historical learned-token
recipe, not native all-patch DINO-WM; model recipes are unmatched.}
\label{tab:shape-control}
\end{table}

Two prespecified one-seed controls are pending for each room count: LpWM with
the spatial 8$\times$8, IDM 3.125, Product-Laplace $\mu=-1,p=1$ recipe used in
the color study, and native all-patch DINO-WM. LpWM tests whether the sparse
global objective is the missing successful family; all-patch DINO-WM tests
whether the historical learned global MLP token, rather than shape information
in the frozen patch representation, caused the weak DINO row. Neither pending
run is used in the table or conclusions.

\section{Provenance, Limitations, and Open Validations}

\paragraph{What the localization does and does not explain.}
The observation is solvable because native LeWM encoder patch tokens and selected
checkpoints from all three comparator families retain hue.  Across the complete LeWM TR2 factorial,
nonlearned 8$\times$8 pooling remains at 100\% hue accuracy, selected global
states span 50.3--92.7\%, and projections span 45.0--56.3\%.  On the
prespecified spatial checkpoint, fixed random 192-D compression reaches 71.1\%,
the learned aggregator 66.0\%, and projection 51.6\%.  These results show a
stagewise decline in hue accessibility along the learned global path but do
not identify the mechanism producing it.  The audited LeWM projector has final
validation-epoch average per-dimension variance .974 and effective rank 45.1 of 192, so the
failure is not scalar variance collapse, although its effective rank is below
the selected PLDM/LpWM controls (72.1/65.8).  Because the current projected
state already has a weak hue readout, the predictor is not the earliest audited
forward stage with weak accessibility; predictor-loss gradients during
training remain unresolved.  Less frequent color-conditioned consequences and the
representation objective remain hypotheses, not identified causes.

\begin{table*}[t]
\centering
\small
\setlength{\tabcolsep}{4pt}
\begin{tabular}{@{}p{.18\textwidth}p{.34\textwidth}p{.36\textwidth}@{}}
\toprule
Explanation & Diagnostic evidence & Status \\
\midrule
Observation lacks hue & LeWM native patches reach 96\%; DINO/PLDM/LpWM retain hue & Rejected for audited images \\
Decoder-only underfit & Ridge and validation-selected MLP also fail at LeWM projection & Insufficient as sole explanation \\
Scalar variance collapse & Projector variance .974 and effective rank 45.1; geometry remains accurate & Rejected as scalar collapse; selective/dimensional loss remains \\
Stagewise global-path decline & Native patches 96--100\%; projected states 45--56\%; fixed random and learned maps differ & Observed; causal mechanism not identified \\
Global readout alone & LeWM CLS/MLP/spatial MAE is 30.69/28.35/25.97; PLDM spatial and LpWM CLS are stronger & LeWM is consistently weak; no readout type is universally sufficient \\
Final projector map & LeWM stays 65.8--68.7\% through GELU, then reaches 51.6\%; PLDM stays at 100\% & Localized at one checkpoint; training cause open \\
Predictor is earliest weak stage & Current projected state is already weak; predicted semantic MAE is similar & Rejected in audited forward order; training gradients remain open \\
Missing temporal alignment & TR2 LeWM is unchanged from weight 0 to 3.125; TR3 spatial improves & Rejected as sole TR2 explanation \\
Less frequent consequence & Teleports occur in 10.21\% of training-split target intervals; TR3 differs in cue support/task statistics & Plausible but not sufficient or identified \\
More distinct colors & Spatial TR3 succeeds, but room count, trajectories, exposure, and updates also change & Open; the prespecified TR4 study is a boundary test, not a cardinality control \\
Regularizer/objective & Matched-architecture/data VISReg comparison yields one 100\% hue readout, not replicated by two seeds or binding & Consistent with sensitivity; causal influence unresolved \\
\bottomrule
\end{tabular}
\caption{Candidate explanations for the audited TR2 LeWM failure.  ``Rejected''
means inconsistent with these diagnostics, not impossible in another setting.}
\label{tab:explanations}
\end{table*}

The audited \lewm{} objective combines prediction and SIGReg (.09)
\citep{balestriero2025lejepa,maes2026leworldmodel} with our IDM extension.
The \pldm{} objective combines prediction with VICReg-inspired VC terms
(variance 18, covariance 12, temporal variance .7, alignment .2; IDM 0)
\citep{bardes2022vicreg,sobal2025latentplanning}.  Our \lpwm{} adaptation uses
prediction, ReLU Product-Laplace InfoReg (10; $\mu=-1,p=1$), and IDM 3.125
\citep{kuang2026rectified}; \dinowm{} predicts frozen pretrained DINOv2 patch
targets.

For the selected TR2 global controls, LeWM, PLDM, and LpWM share the dataset,
seed, 32,550 updates, trainable random ViT-tiny, spatial 8$\times$8 aggregator,
192--2048--192 projector, and six-layer 16-head predictor.  Their checkpoints
have the same 307 tensor-state keys/shapes (20,426,736 elements); decoders see
detached states.  LeWM and PLDM differ principally in representation objective,
while LpWM also rectifies projector and prediction-head outputs.  This makes
the objective the sharpest remaining hypothesis, but coefficients and losses
are not scale-matched, IDM differs for PLDM, and trained initializations are not
paired; it is not a causal ablation.  Exact contracts are in
\path{data/global_model_contracts.csv}.

The 30-run core experiment is frozen by an anonymous source snapshot and
artifact hashes; the private experiment mapping will be released after review.
The follow-up evidence adds 27 aggregation runs (nine conditions over three
seeds) and 28 one-seed temporal-alignment runs.  All 55 finished ten epochs
with final decoder and
location-probe artifacts; their anonymous per-run metrics are frozen in the two
new plotting tables.
Dataset SHA256 values are
\texttt{e134677f...df042} (\trtwo{}) and
\texttt{61c576d...da617} (\trthree{}).  The paper data distribute stable
anonymous experiment IDs and exact plotting inputs for the core tables, plus
media hashes for the decoder audit and checkpoint/split hashes for the
stagewise and counterfactual audits.  All 30 core runs finished ten epochs.
The completion manifest records one absent DINO-WM W\&B energy summary key;
the corresponding final local image is retained and visually valid, so this is
an upload/completion-metadata exception rather than a missing training result.
Figure crops use unscaled pixels from final-epoch training videos; the selection rule and
unmatched DINO frame are disclosed.

The central limitations are that the stagewise, IDM, and temporal-alignment
audits use one training seed; the new aggregation grid has only three seeds;
TR2/TR3 have unequal update counts; model recipes differ in objective and IDM;
the corrected four-color boundary matrix is pending and remains unobserved;
the LpWM and all-patch DINO Shape controls are also pending;
and semantic MAE
composes identity with localization.  Diagnostic decoders may
also underfit, although independent ridge and MLP readouts support the same
LeWM conclusion.  The \trthree{} result explicitly falsifies universal LeWM
color blindness; the fixed-layout pilot is not a hue test.

The stagewise tables preserve checkpoint, dataset, configuration, sidecar, and
episode-split hashes, and the protocol is specified above.  Its one-off runner
was intentionally temporary and is not part of the anonymous artifact.  A
clean checkpoint-to-table evaluator is therefore an artifact-release task, not
an unreported experiment or a source of additional results in this draft.

The original LeWorldModel study also reports weaker performance on simple
Two-Room and hypothesizes that low data diversity or intrinsic dimensionality
can make high-dimensional isotropic-Gaussian SIGReg matching difficult
\citep{maes2026leworldmodel}.  This motivates our TR2/TR3 and regularizer
controls, but neither our unequal-room comparison nor that hypothesis identifies
the mechanism.

VISReg separates a variance-based scale penalty from a sliced-Wasserstein shape
penalty; its authors show that the scale-loss gradient approaches a constant as
embeddings approach collapse \citep{wu2026visreg}.  This property concerns
collapse prevention, not semantic cue retention, so we evaluate VISReg only as
a diagnostic objective-family intervention.  Concurrent VIS-WM work studies
planning rather than color retention and reports VISReg weight 4.5
\citep{wu2026viscore}; our prespecified weight 1.5 experiment is therefore not a
reproduction of that recipe.

The prespecified experiment is a matched-architecture, matched-data recipe
swap, not a regularizer-strength-matched test.  Three seeds each use spatial
8$\times$8 LeWM, IDM 0, ten epochs, and 1024 random directions with either
SIGReg weight .09 or VISReg weight 1.5.  The latter direction count is reduced
from the VISReg paper's $K=4096$ setting.  Method-specific coefficients are not
loss- or gradient-matched: on one frozen epoch-0 state and validation batch,
the weighted regularizer-gradient norm was 9.32$\times$ the prediction gradient
for SIGReg but 2.33$\times$ for VISReg (cosine .675).  Different accelerator
classes under BF16 are another implementation nuisance.  E049's 100\% fixed-slot
hue ridge readout shows that this LeWM architecture can retain accessible TR2
hue under one alternative recipe/seed; it does not show that SIGReg causes cue
loss or that VISReg reliably prevents it.
Before inspecting final outcomes, we designated projector visible-hue balanced
accuracy on the frozen episode-disjoint 525/75/150 split as primary.  Secondary
endpoints are destination-binding accuracy, null-adjusted semantic-hub MAE,
and decoder separation.  We report every completed seed, select no checkpoint
or hyperparameter using test outcomes, and interpret the objective contrast
only if all three seeds and final diagnostics complete.

All six runs and fit-end diagnostics completed.  \Cref{tab:visreg-swap} reports
every seed.  SIGReg remains near chance on hue and binding.  VISReg produces one
perfect hue readout but two near-chance readouts; its binding readouts are all
near chance.  The perfect-hue seed also has a low 4.00-pixel green-hub MAE but a
26.75-pixel orange-hub MAE (not shown), worsened positional mean MAE (15.56
pixels), and 50\% decoder macro color accuracy.  Visual inspection is consistent
with one-color preservation rather than faithful two-color reconstruction, but
the aggregate decoder statistics alone do not identify which color is lost.
This is a heterogeneous one-seed fixed-slot hue result, not a reproducible
semantic-binding result.

\begin{table*}[t]
\centering
\small
\setlength{\tabcolsep}{5pt}
\begin{tabular}{@{}llrrrrrr@{}}
\toprule
ID & Recipe & Seed & Hue BA & Binding BA & Green hub & Position & Decoder sep. \\
 & & & (\%) & (\%) & MAE (px) & MAE (px) & (\%) \\
\midrule
E046 & SIGReg & 3072 & 50.0 & 49.9 & 25.57 & 6.10 & 7.80 \\
E047 & SIGReg & 4072 & 50.8 & 47.8 & 26.10 & 6.35 & 6.04 \\
E048 & SIGReg & 5072 & 47.0 & 47.9 & 25.53 & 6.77 & 1.35 \\
E049 & VISReg & 3072 & \textbf{100.0} & 50.4 & \textbf{4.00} & 15.56 & \textbf{67.22} \\
E050 & VISReg & 4072 & 54.3 & 47.6 & 25.82 & 7.04 & 3.41 \\
E051 & VISReg & 5072 & 50.0 & 56.3 & 26.58 & 7.58 & 1.20 \\
\bottomrule
\end{tabular}
\caption{\textbf{Matched-architecture, matched-data LeWM recipe comparison.}  Hue and binding
are independently validation-selected projector ridge readouts on the same 150
episode test split.  Lower MAE is better.  Decoder separation uses every cue in
the matched final video.  Wilson intervals for all non-perfect raw accuracies
include chance; E049's raw hue interval is [97.5\%, 100\%].  Training seeds are
not uncertainty estimates for the original 12-configuration factorial.}
\label{tab:visreg-swap}
\end{table*}

\paragraph{Causal validations left open.}
{\footnotesize
\begin{enumerate}[leftmargin=*,itemsep=0pt,topsep=2pt]
    \item Compare SIGReg, VISReg, PLDM-style VC, and Product-Laplace InfoReg under a prespecified
    strength-matching rule.
    \item Repeat the same linear and validation-selected MLP readouts on newly
    generated episodes unseen during world-model training.
    \item Match decoder capacity and budget and audit identical
    episode-disjoint frames.
\end{enumerate}
}

\clearpage
\begin{figure*}[p]
    \centering
    \includegraphics[width=\textwidth]{figures/appendix/lewm_factorial_tr2_low.png}
    \caption{\textbf{\trtwo{} decoder factorial, IDM 0--.025.} Columns compare
    native CLS and spatial 8$\times$8.  Spatial decoders remain olive while CLS
    is faint or gray. Within each column the input is left, reconstruction is
    right, and O/D marker crops appear below. Embedded step labels are W\&B
    media indices, not optimizer steps or epochs.}
    \label{fig:factorial-tr2-low}
\end{figure*}

\begin{figure*}[p]
    \centering
    {\small\textbf{\trtwo{} --- native CLS}\hfill
    \textbf{\trtwo{} --- spatial 8$\times$8}\par}
    \includegraphics[width=\textwidth]{figures/appendix/lewm_factorial_tr2_high.png}
    \caption{\textbf{\trtwo{} decoder factorial, IDM .125--3.125.} The matched
    frame rule is identical to \cref{fig:factorial-tr2-low}; averaging persists
    at the larger IDM weights. Embedded step labels are W\&B media indices.}
    \label{fig:factorial-tr2-high}
\end{figure*}

\begin{figure*}[p]
    \centering
    \includegraphics[width=\textwidth]{figures/appendix/lewm_factorial_tr3_low.png}
    \caption{\textbf{\trthree{} decoder factorial, IDM 0--.025.} Spatial
    decoders separate green, orange, and blue; CLS is faint or gray. Within
    each column the input is left, reconstruction is right, and O/D marker
    crops appear below. Embedded step labels are W\&B media indices.}
    \label{fig:factorial-tr3-low}
\end{figure*}

\begin{figure*}[p]
    \centering
    {\small\textbf{\trthree{} --- native CLS}\hfill
    \textbf{\trthree{} --- spatial 8$\times$8}\par}
    \includegraphics[width=\textwidth]{figures/appendix/lewm_factorial_tr3_high.png}
    \caption{\textbf{\trthree{} decoder factorial, IDM .125--3.125.} Spatial
    color separation remains visible but is non-monotonic at high IDM. Embedded
    step labels are W\&B media indices.}
    \label{fig:factorial-tr3-high}
\end{figure*}

\end{document}
